% metodos_simples_vma.tex
% Métodos simples VMA — LoveEarthHacker edition
% Víctor Manzanares Alberola · 2026
% Compilar: pdflatex metodos_simples_vma.tex (×2)

\documentclass[11pt,a4paper,twoside]{book}

\usepackage[utf8]{inputenc}
\usepackage[T1]{fontenc}
\usepackage[spanish,es-tabla]{babel}
\usepackage{amsmath,amssymb,amsfonts}
\usepackage{geometry}
\usepackage{graphicx}
\usepackage{booktabs}
\usepackage{enumitem}
\usepackage{hyperref}
\usepackage{xcolor}
\usepackage{tikz}
\usepackage{listings}
\usetikzlibrary{arrows.meta,positioning,shapes.geometric,calc,fit,backgrounds}

\geometry{margin=2.4cm}
\hypersetup{colorlinks=true, linkcolor=teal!60!black, urlcolor=teal!50!black}

\definecolor{vmaTeal}{RGB}{0,120,110}
\definecolor{vmaGold}{RGB}{200,150,40}
\definecolor{vmaBg}{RGB}{245,248,250}

\newcommand{\LoveEarthHacker}{\textsc{LoveEarthHacker}}
\newcommand{\VMA}{VMA}
\newcommand{\MDC}{\textsc{Mdc}}
\newcommand{\Kthree}{K\textsubscript{3}}

\lstset{
  basicstyle=\ttfamily\small,
  frame=single,
  backgroundcolor=\color{vmaBg},
  breaklines=true,
  columns=fullflexible
}

% ─── Portada ───────────────────────────────────────────────────────────────
\title{%
  {\Huge\bfseries Métodos Simples}\\[0.6em]
  {\Large Herramientas mínimas del ecosistema \VMA}\\[1.2em]
  {\normalsize Implicaciones · Diagramas · Ejecución inmediata}
}
\author{%
  Víctor Manzanares Alberola\\
  \small EPSA · UPV Alcoi · 1477 / 33$\times$1\\[0.4em]
  \small Serie \LoveEarthHacker
}
\date{Julio 2026}

\begin{document}
\frontmatter
\maketitle

\begin{abstract}
Este libro recoge los \textbf{métodos más simples y ejecutables} del corpus
\VMA: aquellos que se pueden explicar en una página, dibujar en un diagrama
y probar en menos de un minuto con \texttt{vma-run} o \texttt{just run}.
No sustituye los papers completos; es la puerta de entrada.
La serie \LoveEarthHacker{} nombra el espíritu: tecnología útil, abierta,
orientada a la Tierra y al rigor, no al postureo.
\end{abstract}

\tableofcontents
\listoffigures
\listoftables

\mainmatter

% ═══════════════════════════════════════════════════════════════════════════
\chapter{Lista de herramientas (mapa rápido)}
% ═══════════════════════════════════════════════════════════════════════════

\begin{table}[htbp]
\centering
\caption{Herramientas simples y comando \texttt{vma-run}}
\label{tab:herramientas}
\begin{tabular}{@{}llll@{}}
\toprule
\textbf{Método} & \textbf{Qué hace} & \textbf{Complejidad} & \textbf{Comando} \\
\midrule
K3 auditoría    & Hash industrial RAM/Flash & $O(n)$ bloques & \texttt{k3 --text "..."} \\
\MDC{} factor   & Factorizar $N$ por salto cinemático & 4 puntos & \texttt{mdc 10403} \\
Criva           & Estimar $\pi(x)/x$ iterativo & $<10$ iter. & \texttt{criva 1000} \\
Criba 6$k{\pm}1$ & Primos en clase modular & $O(n\log\log n)$ & \texttt{criba --limit 500} \\
Fase K=3 XOR    & Amplificador fase oculta & 16 estados & \texttt{phase} \\
Hurto grav.     & Asimetría temporal Quijote & simulación & \texttt{hurto} \\
AntiPC UDP      & Red local como cómputo & 4 hubs & \texttt{antipc} \\
Aleatorovix     & Entropía + criba 6$k{\pm}1$ & C/Python & \texttt{aleatorovix/} \\
Gemelo ZypyZape & Inercia sintética red & Python+plot & \texttt{gemelos/} \\
\bottomrule
\end{tabular}
\end{table}

\begin{figure}[htbp]
\centering
\begin{tikzpicture}[
  node distance=1.1cm and 1.4cm,
  box/.style={draw=vmaTeal, rounded corners, fill=vmaBg, minimum width=2.6cm,
              minimum height=0.9cm, align=center, font=\small},
  arr/.style={-{Stealth[length=2.5mm]}, thick, vmaTeal}
]
  \node[box] (run) {\textbf{RUN.bat}\\\texttt{vma-run.exe}};
  \node[box, below left=of run] (mat) {Matemáticas\\MDC·Criva·Criba};
  \node[box, below=of run] (k3) {Industrial\\K3·Fase K=3};
  \node[box, below right=of run] (eng) {Energía\\Hurto·ZypyZape};
  \node[box, below=1.8cm of run] (net) {Red\\AntiPC·Aleatorovix};
  \draw[arr] (run) -- (mat);
  \draw[arr] (run) -- (k3);
  \draw[arr] (run) -- (eng);
  \draw[arr] (run) -- (net);
  \node[below=0.3cm of net, font=\footnotesize\itshape] {LoveEarthHacker · just run};
\end{tikzpicture}
\caption{Arquitectura mínima: un ejecutable, cuatro familias de métodos.}
\label{fig:mapa}
\end{figure}

% ═══════════════════════════════════════════════════════════════════════════
\chapter{MDC — Método Diofántico Cinemático}
% ═══════════════════════════════════════════════════════════════════════════

\section{La idea en una frase}
En lugar de probar todos los candidatos, mides la \textbf{parte decimal}
$d(t)=\{g(t)\}$ en cuatro puntos consecutivos, calculas velocidad media
$\bar v$ y \textbf{saltas} al $t^\ast$ donde $d(t^\ast)=\delta$.

\section{Fórmula mínima (factorización)}
Para factorizar $N$, define
\[
  d(m) = \operatorname{frac}\!\left(\frac{N}{2(2m+3)}\right), \qquad \delta = \tfrac12.
\]
Si $d(m^\ast)=\tfrac12$, entonces $2m^\ast+3 \mid N$.

\begin{figure}[htbp]
\centering
\begin{tikzpicture}[scale=0.95]
  \draw[->] (0,0) -- (8,0) node[right] {$t$};
  \draw[->] (0,0) -- (0,3) node[above] {$d(t)$};
  \foreach \x/\lab in {1/$t_0$, 2.5/$t_0{+}s$, 4/$t_0{+}2s$, 5.5/$t_0{+}3s$, 7/$t^\ast$}
    \fill[vmaTeal] (\x, {1.2+0.35*sin(\x*60)}) circle (2.5pt);
  \draw[dashed, vmaGold, thick] (0,1.5) -- (8,1.5) node[right] {$\delta$};
  \draw[thick, vmaTeal] plot[smooth] coordinates {(0.5,0.8)(2,1.4)(3.5,0.6)(5,1.8)(6.5,0.4)(7.5,1.5)};
  \node[font=\small] at (4,-0.6) {Diente de sierra $\rightarrow$ salto cinemático};
\end{tikzpicture}
\caption{MDC: cuatro muestras, una predicción. Sin barrido lineal.}
\end{figure}

\section{Implicaciones}
\begin{itemize}[noitemsep]
  \item \textbf{Matemáticas:} unifica factorización y primos de Wieferich (mismo motor, distinto $\delta$).
  \item \textbf{33$\times$1:} núcleo del factorizador territorial; software puro.
  \item \textbf{Límite:} candidato primo $\neq$ demostración; validar con prueba clásica.
\end{itemize}

\section{Probar}
\begin{lstlisting}
vma-run mdc 10403    # 101 x 103
vma-run mdc 15251    # 101 x 151
\end{lstlisting}

% ═══════════════════════════════════════════════════════════════════════════
\chapter{Criva — densidad racional de primos}
% ═══════════════════════════════════════════════════════════════════════════

\section{La idea en una frase}
Estimas $\pi(x)/x$ mezclando el producto de Euler (cribas por capas) con
contracción iterativa hacia $1/\log x$:
\[
  D_{n+1} = \frac{D_n + T_n}{2}, \qquad T_n = D_n\Bigl(1-\frac{1}{\log(x+1)}\Bigr).
\]

\begin{figure}[htbp]
\centering
\begin{tikzpicture}[
  layer/.style={draw, fill=vmaBg, minimum height=0.55cm, minimum width=3.2cm},
  font=\small
]
  \node[layer] (l0) at (0,0) {Capa 0: todos los enteros};
  \node[layer, below=0.15cm of l0] (l1) {Capa 1: sin múltiplos de 2};
  \node[layer, below=0.15cm of l1] (l2) {Capa 2: sin 2, 3};
  \node[layer, below=0.15cm of l2] (l3) {Capa 3: sin 2, 3, 5};
  \node[below=0.2cm of l3, font=\footnotesize] {$\cdots$ convergencia $\rightarrow\ \pi(x)/x$};
  \draw[->, thick, vmaTeal] (-2.2,0.3) -- (-2.2,-2.2) node[midway, left, align=right] {Criva\\itera};
\end{tikzpicture}
\caption{Criva como capas de exclusión + refinamiento.}
\end{figure}

\section{Implicaciones}
\begin{itemize}[noitemsep]
  \item Publicable como estimador ligero (sin \texttt{sympy} obligatorio).
  \item Complementa MDC: uno localiza factores, otro describe densidad global.
  \item Error $<0.1\%$ para $x>100$ en pocas iteraciones (modo conservador).
\end{itemize}

% ═══════════════════════════════════════════════════════════════════════════
\chapter{Criba modular 6\texorpdfstring{$k\pm1$}{k±1}}
% ═══════════════════════════════════════════════════════════════════════════

\section{La idea en una frase}
Solo candidatos $6k\pm1$; para cada primo $p$ saltas $+2p$ y $+4p$
(tocas ambas clases sin barrer todos los enteros).

\section{Implicaciones}
\begin{itemize}[noitemsep]
  \item Base de Aleatorovix y Libro~5; memoria $\sim 4/9$ del tamiz clásico.
  \item Conecta con MDC: espacio de búsqueda $L_1=\{6k\pm1\}$.
\end{itemize}

% ═══════════════════════════════════════════════════════════════════════════
\chapter{\Kthree{} — auditoría y amplificador de fase}
% ═══════════════════════════════════════════════════════════════════════════

\section{K3 hash (industrial)}
Motor de auditoría \textbf{no criptográfico}: integridad de memoria, telemetría,
fingerprinting. Tres modos: usual, propio-A (stride), propio-B (marcas).

\begin{figure}[htbp]
\centering
\begin{tikzpicture}[
  block/.style={draw, rounded corners, fill=white, minimum width=1.8cm, minimum height=0.7cm, font=\footnotesize}
]
  \node[block] (in) {Entrada};
  \node[block, right=0.5cm of in] (k3) {Compresión\\\Kthree{}};
  \node[block, right=0.5cm of k3] (out) {Hash 32b};
  \draw[->, thick] (in) -- (k3) -- (out);
  \node[below=0.5cm of k3, font=\footnotesize\itshape] {firma anti-tamper en bloque 0};
\end{tikzpicture}
\caption{Flujo K3: un paso, salida determinista.}
\end{figure}

\section{Amplificador de fase K=3 XOR}
Dos estados difieren en fase oculta $f_1$; tras un ciclo Toffoli+XOR+shift,
la diferencia aparece en $v_2$ \textbf{para siempre}.

\section{Implicaciones}
\begin{itemize}[noitemsep]
  \item K3: producto B2B inmediato (\texttt{vma-k3}, TRL 3--4).
  \item Fase K=3: detección de paridad de fase; puente a ciberseguridad geométrica.
\end{itemize}

% ═══════════════════════════════════════════════════════════════════════════
\chapter{Hurto gravitatorio y Quijote (versión simple)}
% ═══════════════════════════════════════════════════════════════════════════

\section{La idea en una frase}
Rompes simetría \textbf{temporal}: extraes energía en bajada rápida,
reinvertís menos en subida lenta. El ciclo cerrado conserva energía total,
pero el \textbf{timing} del robo de fase produce ganancia neta en simulación.

\begin{figure}[htbp]
\centering
\begin{tikzpicture}[scale=0.9]
  \draw[thick, vmaTeal] (0,0) sin (45:2.2) cos (45:2.2);
  \draw[thick, vmaTeal] (0,0) sin (135:1.2) cos (135:1.2);
  \node[font=\small] at (1.5,1.5) {bajada rápida};
  \node[font=\small] at (-0.9,0.5) {subida lenta};
  \draw[->, vmaGold, thick] (2.5,0) arc (0:300:2.5);
  \node[right, font=\footnotesize] at (2.6,-0.2) {$W_{\mathrm{neto}}>0$ en modos PAPER/SINUS};
\end{tikzpicture}
\caption{Asimetría temporal: no es perpetuum mobile, es sincronización de fase.}
\end{figure}

\begin{table}[htbp]
\centering
\caption{Resultados simulación NREL 5\,MW (90\,s) — \texttt{quijote\_results.csv}}
\begin{tabular}{@{}ccrrr@{}}
\toprule
$N$ pales & Modo & $\bar P_{\mathrm{net}}$ [kW] & $\eta_{\mathrm{hurto}}$ & $\Delta P_{\mathrm{grid}}$ [\%] \\
\midrule
3 & PAPER & $+9.0$  & $3.0\times$ & $4.16$ \\
7 & PAPER & $+31.1$ & $3.7\times$ & $4.31$ \\
3 & CONTROL & $-15.4$ & --- & --- \\
\bottomrule
\end{tabular}
\end{table}

\section{Implicaciones}
\begin{itemize}[noitemsep]
  \item \textbf{Seguro hoy:} inercia sintética ZypyZape (firmware, sin hardware nuevo).
  \item \textbf{Pendiente:} modo CONTROL y prototipo físico (IEC).
  \item \textbf{33$\times$1:} argumento energético territorial.
\end{itemize}

% ═══════════════════════════════════════════════════════════════════════════
\chapter{AntiPC y Aleatorovix (red simple)}
% ═══════════════════════════════════════════════════════════════════════════

\section{AntiPC — red como cómputo}
Los hubs UDP procesan fragmentos; el maestro solo fusiona bits en búfer
lock-free. Sin \texttt{if/else} en el camino crítico (Grafcet algebraico).

\begin{figure}[htbp]
\centering
\begin{tikzpicture}[
  hub/.style={draw, circle, fill=vmaBg, minimum size=1cm, font=\footnotesize},
  m/.style={draw, rounded corners, fill=white, minimum width=2cm, minimum height=0.8cm, font=\small}
]
  \node[hub] (h1) at (0,1.5) {H1};
  \node[hub] (h2) at (-1.5,-0.5) {H2};
  \node[hub] (h3) at (1.5,-0.5) {H3};
  \node[hub] (h4) at (0,-2) {H4};
  \node[m] (master) at (4,0) {Maestro\\búfer RAM};
  \foreach \h in {h1,h2,h3,h4}
    \draw[->, thick, vmaTeal] (\h) -- (master);
  \node[below=0.4cm of h4, font=\footnotesize] {$\sim$2000 pkt/s en localhost (4 hubs)};
\end{tikzpicture}
\caption{AntiPC: delegar, fusionar, no serializar en el OS.}
\end{figure}

\section{Aleatorovix — entropía + criba}
Organismo que roba bits de pila/heap/nanosegundos, decide salto $6k\pm1$
y borra su rastro (\emph{criba desmemoriada}). Concepto Libro~5; C en
\texttt{organismo\_lila\_v99.c}.

\section{Implicaciones}
\begin{itemize}[noitemsep]
  \item Hardware barato (hubs $\sim$60--80\,€/nodo): documento \texttt{cpu-antipc}.
  \item Speedup real requiere benchmark frente a RDMA/DPDK --- honestidad técnica.
  \item \LoveEarthHacker: computación distribuida consciente, no minería extractiva.
\end{itemize}

% ═══════════════════════════════════════════════════════════════════════════
\chapter{LoveEarthHacker — marco y 33\texorpdfstring{$\times$}{x}1}
% ═══════════════════════════════════════════════════════════════════════════

\section{Qué significa \LoveEarthHacker{} aquí}
No es un producto cerrado: es el \textbf{contrato ético-técnico} del corpus:
\begin{enumerate}[noitemsep]
  \item Métodos \textbf{simples} primero (ejecutables, reproducibles).
  \item Implicaciones declaradas (qué está probado y qué es simulación).
  \item Tierra y soberanía: energía, red, mates --- no dependencia de monopolio.
  \item 33$\times$1: 33 años cooperación $\leftrightarrow$ 1 paquete tecnológico territorial.
\end{enumerate}

\begin{figure}[htbp]
\centering
\begin{tikzpicture}
  \node[draw, rounded corners, fill=vmaBg, inner sep=12pt, align=center, text width=10cm] {
    \textbf{LoveEarthHacker}\\[0.3em]
    Método simple $\rightarrow$ Diagrama $\rightarrow$ \texttt{vma-run demo}\\
    $\rightarrow$ Implicación (industrial / energía / red / 33$\times$1)
  };
\end{tikzpicture}
\caption{Ciclo de validación mínima de cada herramienta.}
\end{figure}

\section{Orden de lectura recomendado}
\begin{enumerate}
  \item Este libro (mapa).
  \item \texttt{just run/LEEME.txt} (ejecución).
  \item \texttt{00\_PROMPT\_DEFINITIVO\_33x1.md} (marco global).
  \item Papers: Wieferich, hurto, teorema K=3 XOR.
\end{enumerate}

% ═══════════════════════════════════════════════════════════════════════════
\appendix
\chapter{Comandos \texttt{vma-run} (referencia)}
% ═══════════════════════════════════════════════════════════════════════════

\begin{lstlisting}
RUN.bat                              # demo completa
bin\vma-run.exe demo
vma-run mdc 10403
vma-run k3 --text "LoveEarthHacker"
vma-run criva --compare 100,1000,10000
vma-run criba --limit 1000 --show 25
vma-run phase
vma-run hurto
vma-run antipc --hubs 4 --packets 500
\end{lstlisting}

\chapter{Rutas en \texttt{just run}}
\begin{itemize}
  \item \texttt{vma-run/} --- ejecutable unificado
  \item \texttt{vma-k3/} --- motor K3
  \item \texttt{aleatorovix/} --- Lila + MDC v6 demo
  \item \texttt{cpu-antipc/} --- especificación hubs + \texttt{.cpp}
  \item \texttt{gemelos/} --- Quijote / ZypyZape
  \item \texttt{hurto-gravitatorio/} --- CSV + reglas paper
\end{itemize}

\backmatter
\chapter*{Colofón}
\addcontentsline{toc}{chapter}{Colofón}
\LoveEarthHacker{} · Métodos simples VMA · Víctor Manzanares Alberola · 2026\\
Compilar: \texttt{pdflatex metodos\_simples\_vma.tex} (dos pasadas).\\
Repositorio canónico: \texttt{just run} en Desktop.

\end{document}